\section{Reproducing the verification}\label{app:reproduce}

Everything needed to repeat \cref{sec:verification} is in the repository
\begin{center}
\url{https://github.com/vincitamore/riemann-lean-account}
\end{center}
together with the source of this paper, the reading notes, the delegated reports and the audit.
The archive itself is not redistributed there; it is fetched from the author's site and checked
against the hash of \cref{sec:verification}.

\subsection*{Requirements}
A Linux machine with at least 32 cores, 128\,GB of memory and 150\,GB of free disk gives the
times quoted here. The scripts in \lean{verify/} assume \lean{bash} with \lean{git},
\lean{curl}, \lean{zstd}, \lean{tar} and \lean{sha256sum}, and \lean{python3} for the scan;
\lean{setup.sh} installs \lean{elan} if it is absent.

\subsection*{Steps}
\begin{enumerate}
\item \lean{setup.sh} downloads the fourteen parts, concatenates them, checks the SHA-256,
  unpacks, installs the pinned toolchain and fetches the Mathlib build cache for the pinned
  commit.
\item \lean{scan.py} is the static inspection of \cref{sec:verification}: it reads the
  downloaded file without unpacking it and needs nothing else.
\item \lean{build.sh} runs \lean{lake build Main} with every log line time-stamped and records
  the exit code.
\item \lean{janitor.sh} must run beside the build. Lake writes, for every module, a file listing
  the module's whole import closure; for this tree those files would total about 900\,GB. The
  script deletes each one once its module's compiled file exists, which Lake tolerates. With it
  the disk peaks near 110\,GB.
\item \lean{check.sh} prints the statement of the main theorem, the module in which Mathlib's
  statement is declared, and the axioms of the five theorems.
\item \lean{cut.sh} builds and runs the program \lean{cut/CutMain.lean}, which computes the
  dependency cone of \cref{sec:cone} and writes the list of every declaration of the archive
  inside it and outside it.
\item \lean{replay.sh} builds and runs the program \lean{replay/ReplayMain.lean}, the kernel
  replay of \cref{sec:verification}. Our run used 28 threads, set with
  \lean{LEAN_NUM_THREADS}.
\end{enumerate}
The outputs of our own run are in \lean{verify/results/}, with an index that says how each
file was produced: the scan, the head and tail of the build log, the output of the check, the
cone summary with the complete lists of declarations inside and outside the cone (compressed),
the replay summary with the lists of axioms and of declarations not replayed, and the output
of the corollaries file.

\subsection*{Cost}
The build took \buildhours{} hours of wall time on 32 cores, the cone about ten minutes and the
replay \replayminutes{} minutes. On rented hardware the whole verification cost under thirty US
dollars.
